\documentclass{article}
\usepackage[english,russian]{babel}
\usepackage{textcomp}
\usepackage{geometry}
  \geometry{left=2cm}
  \geometry{right=1.5cm}
  \geometry{top=1.5cm}
  \geometry{bottom=2cm}
\usepackage{tikz}
\usepackage{multicol}
\usepackage{hyperref}
\usepackage{listings}
\pagenumbering{gobble}

\lstdefinestyle{csMiptCppStyle}{
  language=C++,
  basicstyle=\linespread{1.1}\ttfamily,
  columns=fixed,
  fontadjust=true,
  basewidth=0.5em,
  keywordstyle=\color{blue}\bfseries,
  commentstyle=\color{gray},
  texcl=true,
  stringstyle=\ttfamily\color{orange!50!black},
  showstringspaces=false,
  numbersep=5pt,
  numberstyle=\tiny\color{black},
  numberfirstline=true,
  stepnumber=1,      
  numbersep=10pt,
  backgroundcolor=\color{white},
  showstringspaces=false,
  captionpos=b,
  breaklines=true
  breakatwhitespace=true,
  xleftmargin=.2in,
  extendedchars=\true,
  keepspaces = true,
  tabsize=4,
  upquote=true,
}


\lstdefinestyle{csMiptCppLinesStyle}{
  style=csMiptCppStyle,
  frame=lines,
}

\lstdefinestyle{csMiptCppBorderStyle}{
  style=csMiptCppStyle,
  framexleftmargin=5mm, 
  frame=shadowbox, 
  rulesepcolor=\color{gray}
}


\lstdefinestyle{csMiptBash}{
breaklines=true,
frame=tb,
language=bash,
breakatwhitespace=true,
alsoletter={*()"'0123456789.},
alsoother={\{\=\}},
basicstyle={\ttfamily},
keywordstyle={\bfseries},
literate={{=}{{{=}}}1},
prebreak={\textbackslash},
sensitive=true,
stepnumber=1,
tabsize=4,
morekeywords={echo, function},
otherkeywords={-, \{, \}},
literate={\$\{}{{{{\bfseries{}\$\{}}}}2,
upquote=true,
frame=none
}




\lstset{style=csMiptCppLinesStyle}
\lstset{literate={~}{{\raisebox{0.5ex}{\texttildelow}}}{1}}


\renewcommand{\thesection}{\arabic{section}}
\makeatletter
\def\@seccntformat#1{\@ifundefined{#1@cntformat}%
   {\csname the#1\endcsname\quad}%    default
   {\csname #1@cntformat\endcsname}}% enable individual control
\newcommand\section@cntformat{Часть \thesection:\space}
\makeatother




\begin{document}
\title{Семинар \#1b: Линковка \vspace{-5ex}}\date{}\maketitle

\section{Преобразование кода на языке C в бинарный код}

\section{Формат файла ELF}

\subsection*{Утилита \texttt{readelf}}
\begin{itemize}
\item Показать ELF заголовок:
\begin{lstlisting}[style=csMiptBash]
readelf -h имя-файла
\end{lstlisting}

\item Показать заголовки секций:
\begin{lstlisting}[style=csMiptBash]
readelf -S имя-файла
readelf -t имя-файла
\end{lstlisting}

\item Показать заголовки сегментов:
\begin{lstlisting}[style=csMiptBash]
readelf -l имя-файла
\end{lstlisting}

\item Показать таблицу символов:
\begin{lstlisting}[style=csMiptBash]
readelf -s имя-файла
\end{lstlisting}

\item Показать динамические секции:
\begin{lstlisting}[style=csMiptBash]
readelf -d имя-файла
\end{lstlisting}

\item Показать таблицу релокаций:
\begin{lstlisting}[style=csMiptBash]
readelf -r имя-файла
\end{lstlisting}

\item Hex-dump секции:
\begin{lstlisting}[style=csMiptBash]
readelf -x имя-секции имя-файла
\end{lstlisting}

\item Показать всё:
\begin{lstlisting}[style=csMiptBash]
readelf -a имя-файла
\end{lstlisting}
\end{itemize}


\subsection*{Утилита \texttt{objdump}}
\begin{itemize}
\item Дизассемблировать секцию .text:
\begin{lstlisting}[style=csMiptBash]
objdump -d имя-файла
objdump -S имя-файла
\end{lstlisting}
\end{itemize}

\section{Процесс статической линковки}

\begin{itemize}
\item \textbf{Парсинг объектных файлов и статичских библиотек}\\
Извлекает все необходимые секции, в частности \texttt{.symtab}. Для статических библиотек просматривает символьный индекс библиотеки и определяет какие объектные файлы из библиотеки нужны, берёт только их.
\item \textbf{Разрешение символов}\\
Строит глобальную таблицу символов
Для каждого неопределённого символа ищет определение в других объектных файлах.
Если символ не найден, то ошибка \texttt{undefined reference}.

\item \textbf{Слияние секций}\\

\item \textbf{Релокация}\\

\item \textbf{Оптимизации и генерация итогового исполняемого файла}\\

\end{itemize}


\section{Процесс динамической линковки}
\begin{itemize}
\item \textbf{Парсинг объектных файлов и статичских библиотек}\\
Извлекает все необходимые секции, в частности \texttt{.symtab}. Для статических библиотек просматривает символьный индекс библиотеки и определяет какие объектные файлы из библиотеки нужны, берёт только их.

\item \textbf{Разрешение символов}\\
Строит глобальную таблицу символов
Для каждого неопределённого символа ищет определение в других объектных файлах.
Если символ не найден, то ошибка \texttt{undefined reference}.

\item \textbf{Слияние секций}\\

\item \textbf{Релокация}\\

\item \textbf{Оптимизации и генерация итогового исполняемого файла}\\

\end{itemize}


\section{Особенности компиляции и линковки в Windows}

\end{document}
